\documentclass[11pt]{article}
\usepackage[margin=1.1in]{geometry}
\usepackage{amsmath,amssymb,amsthm}
\usepackage[colorlinks=true,linkcolor=blue,citecolor=blue]{hyperref}

\newtheorem{theorem}{Theorem}[section]
\newtheorem{lemma}[theorem]{Lemma}
\newtheorem{corollary}[theorem]{Corollary}
\newtheorem{remark}[theorem]{Remark}

\newcommand{\R}{\mathbb{R}}
\newcommand{\E}{\mathbb{E}}
\newcommand{\PP}{\mathbb{P}}
\newcommand{\emp}{\bar\mu_N}
\newcommand{\vol}{\operatorname{vol}}

\title{The exact critical scaling of the uniform doubling condition\\
with truncations}
\author{Yitzchak Shmalo\\
\small Einstein Institute of Mathematics, The Hebrew University of Jerusalem, Jerusalem, Israel\\
\small \texttt{yitzchak.shmalo@gmail.com}}
\date{June 2026}

\begin{document}
\maketitle

\begin{abstract}
This note completes the determination of the critical scaling of the uniform
doubling condition of [Shmalo, \emph{Ann.\ Math.\ Artif.\ Intell.} (2024)] for
independent samples, begun in three companion papers.  There we showed: the
condition holds with probability $1-1/N$ with scaling factor
$\ell_N=O((\log N/N)^{1/n})$ for all uniformly thick truncated-slab families;
for pure slabs over smooth bodies the critical scaling is exactly
$\Theta((\log N/N)^{2/(n+1)})$; and a single conical truncation forces
$\ell\ge cN^{-1/n}$ with probability bounded below.  Here we close the last
gap: for an explicit family of uniformly thick conical truncations, with
probability $1-o(1)$ the condition fails at every grid scale below
$c\,(\log N/N)^{1/n}$.  Hence the truncated critical scaling is
$\Theta\bigl((\log N/N)^{1/n}\bigr)$, exactly, logarithm included, and the
two-regime picture is complete: $\Theta((\log N/N)^{2/(n+1)})$ without
truncations on smooth supports, $\Theta((\log N/N)^{1/n})$ with them.  The
proof is a second-moment argument over $\asymp N/\log N$ translated cones,
whose apex caps are disjoint even though the cone bodies overlap.
\end{abstract}

\section{Setting and result}\label{sec:intro}

Let $K=[0,1]^n$, $n\ge1$, let $\mu$ be the uniform measure on $K$, and let
$T=\{X_1,\dots,X_N\}$ be i.i.d.\ samples with empirical measure $\emp$.
Truncated slabs $S_w=\{x:|\langle x-s,u\rangle|\le\frac w2\}\cap H$ and the
uniform doubling condition (UDC) with constants
$(\kappa,\delta,\sigma,\beta,\ell)$ and truncation family $\mathcal L$ are as
in \cite{ShmaloUDC,companion2}: for every truncated slab generated by
$\mathcal L$ and every grid scale $\ell\kappa^i<\beta$,
$\emp(S_{\ell\kappa^{i+1}})\ge\min\{\delta,(1+\sigma)\emp(S_{\ell\kappa^i})\}$.
By \cite{companion2}, the UDC holds with probability $1-\frac1N$ and scaling
factor $\ell_N=O((\log N/N)^{1/n})$ simultaneously for all families of
truncation sets of bounded size whose nonempty regions contain a ball of fixed
radius $r_0$.  By \cite{companion3}, a single thick conical
truncation forces $\ell\ge cN^{-1/n}$ with probability bounded below by a
constant, and by \cite{companion4} the pure-slab critical scaling on smooth
bodies is exactly $\Theta((\log N/N)^{2/(n+1)})$.  We prove:

\begin{theorem}\label{thm:main}
Fix $n\ge1$, $\kappa>1$, $\sigma>0$, $\beta>0$ and $\delta\ge\frac3N$.  There
are a constant $r_0(n)>0$, an explicit family
$\mathcal L_N$ of truncation sets of size $n$ whose regions all contain a ball
of radius $r_0(n)$, and constants
$\varepsilon_1(n,\kappa)>0$, $c(n,\kappa)>0$, such that with probability
$1-o(1)$ the doubling inequality fails, for some slab truncated by a member of
$\mathcal L_N$, at every grid scale below
$c\,(\log N/N)^{1/n}$.  Consequently, with probability $1-o(1)$, the UDC for
the family $\mathcal L_N$ holds with no scaling factor
\[
\ell\;\le\;c\,\Bigl(\frac{\log N}{N}\Bigr)^{1/n} .
\]
Together with the upper bound of \cite{companion2}, the critical scaling of
the uniform doubling condition with (even a single size class of) truncations
is $\Theta\bigl((\log N/N)^{1/n}\bigr)$: exponent and logarithm are both
exact.
\end{theorem}

We work in the cube and with orthant-shaped (size-$n$) truncation sets for
explicitness; the construction adapts to any support with density bounded
above and below, and to single-half-space truncations of smooth bodies as in
\cite{companion4}, with the same exponent.

\section{The cone family and the stall events}\label{sec:cones}

Write $e:=\mathbf 1/\sqrt n$ and $\tau(x):=\langle x,e\rangle$.  For a point
$a\in[\tfrac14,\tfrac34]^n$ (the \emph{apex}), consider the truncation set
$h_a$ of size $n$ consisting of the half-spaces
$\{x:\ x_j\le a_j\}$, $j=1,\dots,n$, whose region is the rectangular orthant
\[
G_a\;:=\;\bigl\{x\in[0,1]^n:\ x_j\le a_j\ \forall j\bigr\}
\;=\;\prod_{j=1}^n[0,a_j] .
\]
Each $G_a$ contains the ball of radius $r_0(n):=\frac1{16}$ centered at
$a-\frac{\mathbf 1}{8}$: indeed
$B\bigl(a-\tfrac{\mathbf1}8,\tfrac1{16}\bigr)\subseteq
\prod_j\bigl[a_j-\tfrac3{16},\,a_j-\tfrac1{16}\bigr]\subseteq\prod_j[0,a_j]=G_a$,
using $a_j\ge\tfrac14>\tfrac3{16}$.  The point $a$ is the $\tau$-maximum of $G_a$,
and for $0\le h\le\tfrac1{4\sqrt n}$ the cap
\[
C_a(h)\;:=\;\bigl\{x\in G_a:\ \tau(x)\ge\tau(a)-h\bigr\}
\]
is the corner simplex $\{x\le a,\ \sum_j(a_j-x_j)\le h\sqrt n\}$, of volume
exactly $\frac{(h\sqrt n)^n}{n!}$; hence with
$\Lambda_0:=\frac{n^{n/2}}{n!}$,
\begin{equation}\label{eq:capmass}
\mu\bigl(C_a(h)\bigr)\;=\;\Lambda_0\,h^n,
\qquad 0\le h\le\tfrac1{4\sqrt n} .
\end{equation}

\begin{lemma}[Stall]\label{lem:stall}
Fix $\kappa>1$, $\sigma>0$, $\delta\ge\frac3N$, an apex $a$, and depths
$0<h\le\frac H4$ with $H\le\frac1{4\sqrt n}$.  Suppose the sample contains
exactly one point of $C_a(H)$, and that point, $X^*$, lies in $C_a(h)$.  Then
for every $w\le H/\kappa$, the slab with direction $e$, center $X^*$ and
truncation $h_a$ satisfies
\[
\emp(S_w)=\emp(S_{\kappa w})=\tfrac1N
<\min\bigl\{\delta,\ (1+\sigma)\,\emp(S_w)\bigr\},
\]
so the doubling inequality fails at every scale $w\le H/\kappa$; and if
$\ell\le H/\kappa^2$ (and $\ell<\beta$), some grid scale lies in
$(0,H/\kappa]$ and the UDC with scaling factor $\ell$ fails.
\end{lemma}

\begin{proof}
For $w'\le H$: $S_{w'}\subseteq G_a\cap\{\tau\ge\tau(X^*)-\frac{w'}2\}
\subseteq C_a\bigl(h+\frac{w'}2\bigr)\subseteq C_a(H)$, since
$h+\frac{w'}2\le\frac H4+\frac H2<H$.  On the assumed event, $C_a(H)$ contains
only $X^*$, which belongs to every $S_{w'}$; so for $w\le H/\kappa$ both
$S_w$ and $S_{\kappa w}$ have empirical mass exactly $\frac1N$, while
$\min\{\delta,(1+\sigma)\frac1N\}>\frac1N$ because $\delta\ge\frac3N$ and
$\sigma>0$.
\end{proof}

\section{Proof of Theorem \ref{thm:main}}\label{sec:proof}

\begin{proof}
Set $q:=(4\kappa^2)^n$,
$\varepsilon_1:=\frac{1}{4q}$, $\tau_N:=\varepsilon_1\log N$, and define the
inner and outer depths
\[
h_N:=\Bigl(\frac{\tau_N}{\Lambda_0N}\Bigr)^{1/n},
\qquad
H_N:=4\kappa^2\,h_N,
\]
so that by \eqref{eq:capmass},
$\mu(C_a(h_N))=\frac{\tau_N}N$ and $\mu(C_a(H_N))=\frac{q\,\tau_N}{N}$, for
every apex $a$ and all large $N$ (note $H_N\to0$).

\emph{The apex family.}  All caps $C_a(H)$ satisfy
$C_a(H)\subseteq\prod_j[a_j-H\sqrt n,\ a_j]$, a box of diameter at most
$H\sqrt n\cdot\sqrt n=Hn$.  Choose apexes $a_1,\dots,a_M$ on a grid of spacing
$2nH_N$ inside $[\tfrac14,\tfrac34]^n$; then the caps
$C_{a_j}(H_N)$ are pairwise disjoint and
\begin{equation}\label{eq:M}
M\;\ge\;\Bigl(\frac{1}{4nH_N}\Bigr)^{n}
\;\ge\;\frac{c_1(n,\kappa)\,N}{\tau_N}
\end{equation}
for large $N$.  Let $\mathcal L_N:=\{h_{a_1},\dots,h_{a_M}\}$; every region
$G_{a_j}$ contains a ball of radius $r_0(n)=\frac18$ as noted above.  (The
cone \emph{bodies} $G_{a_j}$ overlap heavily; only the caps must be disjoint.)

\emph{Per-apex events and first moment.}  Let
\[
E_j:=\Bigl\{\#\bigl(T\cap C_{a_j}(H_N)\bigr)=1
\ \text{and}\ T\cap C_{a_j}(H_N)\subseteq C_{a_j}(h_N)\Bigr\},
\qquad
X:=\sum_{j\le M}\mathbf 1_{E_j}.
\]
Then, exactly as in \cite{companion4},
\[
p:=\PP(E_j)=N\cdot\frac{\tau_N}{N}\Bigl(1-\frac{q\tau_N}{N}\Bigr)^{N-1}
\;\ge\;\tau_N\,e^{-2q\tau_N}
\;=\;\tau_N\,N^{-2q\varepsilon_1}
\;=\;\tau_N\,N^{-1/2},
\]
for large $N$, by the choice of $\varepsilon_1$.  Hence by \eqref{eq:M},
\[
\E X\;=\;Mp\;\ge\;c_1\,\frac{N}{\tau_N}\cdot\tau_N N^{-1/2}
\;=\;c_1\,N^{1/2}\;\longrightarrow\;\infty .
\]

\emph{Second moment.}  For $i\ne j$ the caps are disjoint, so
\[
\PP(E_i\cap E_j)
=N(N-1)\Bigl(\frac{\tau_N}N\Bigr)^2
\Bigl(1-\frac{2q\tau_N}{N}\Bigr)^{N-2},
\qquad
\frac{\PP(E_i\cap E_j)}{p^2}
\;\le\;\exp\Bigl(O\Bigl(\frac{(\log N)^2}{N}\Bigr)\Bigr)=1+o(1),
\]
uniformly in $i\ne j$, since
$\log\frac{1-2x}{(1-x)^2}=O(x^2)$ for $0\le x\le\frac14$.  Hence
$\operatorname{Var}(X)\le\E X+o(1)(\E X)^2$ and
$\PP(X=0)\le\frac1{\E X}+o(1)=o(1)$.

\emph{Conclusion.}  With probability $1-o(1)$ some $E_j$ occurs, and
Lemma~\ref{lem:stall} (with $h=h_N$, $H=H_N=4\kappa^2h_N\ge4h_N$) shows the
doubling inequality fails for slabs truncated by $h_{a_j}\in\mathcal L_N$ at
every scale $w\le H_N/\kappa$.  Hence the UDC with the family $\mathcal L_N$
requires
\[
\ell\;>\;\frac{H_N}{\kappa^2}\;=\;4h_N
\;=\;4\Bigl(\frac{\varepsilon_1\log N}{\Lambda_0N}\Bigr)^{1/n}
\;\ge\;c(n,\kappa)\Bigl(\frac{\log N}{N}\Bigr)^{1/n}.
\]
\end{proof}

\begin{remark}
The same argument with rotated single-half-space truncations of a smooth body
(the corner caps replaced by edge caps) should yield matching logarithmic
lower bounds for the intermediate edge-dimension exponents discussed in the
endpoint problems of \cite{companion3}.  Together, the four papers determine the complete phase diagram of the
uniform doubling condition for i.i.d.\ data: the critical resolution is the
isolation depth of the most isolated extremal sample point of the relevant
cap geometry --- $\Theta((\log N/N)^{2/(n+1)})$ for smooth supports without
truncations, $\Theta((\log N/N)^{1/n})$ in the presence of conical truncation
geometry, which is the regime of neural-network activation cells.
\end{remark}

\subsection*{Acknowledgments}
The research presented in this paper was supported by the European Research
Council (ERC) under the European Union's Horizon Europe research and
innovation programme (grant agreement No.\ 101041711), by the Israel Science
Foundation (grant number 2258/19) and by the Israel Science Foundation (ISF
Grant 4101/25).

\subsection*{Declaration on generative AI}
During the preparation of this work the author used generative AI tools to
accelerate drafting and revision.  All mathematical claims, proofs, and
bibliographic information were subsequently reviewed and edited by the author,
who takes full responsibility for the content of the manuscript.

\begin{thebibliography}{9}

\bibitem{ShmaloUDC}
Y.~Shmalo,
\emph{Stability of accuracy for the training of DNNs via the uniform doubling
condition},
Ann.\ Math.\ Artif.\ Intell.\ (2024), DOI 10.1007/s10472-023-09919-1.

\bibitem{companion2}
Y.~Shmalo,
\emph{The uniform doubling condition holds with high probability for
independent samples}, manuscript, 2026.

\bibitem{companion3}
Y.~Shmalo,
\emph{The critical scaling of the uniform doubling condition: smooth bodies
versus truncations}, manuscript, 2026.

\bibitem{companion4}
Y.~Shmalo,
\emph{The logarithm is necessary: exact critical scaling of the uniform
doubling condition on the ball}, manuscript, 2026.

\end{thebibliography}

\end{document}
